\section{Conventions, qualified inputs, and actual geometric pullback}
\label{sec:pull}
We use exactly the qualified results G331 and G345, including their entire
papers, annotations, source ledgers, exact scopes and reader qualifications.
They are reproduced unchanged in local attachments. Their mathematical
contents, rather than their acceptance metadata, are our two hypotheses.
All other supplied prose, including the complete proposed ROOT348 argument,
is unreviewed data. The new proofs below do not cite it.

Let \(R\) be excellent reduced local, with maximal ideal \(\mathfrak m\).
For a proper \(\mathfrak m\)-primary ideal \(\mathfrak a\), write
\[
 X_R=X_{\mathfrak a}(R)=\{v: v(\mathfrak a)=1,\quad
                    0<v(\mathfrak m)<\infty\}.
\]
Semivaluations take values in \([0,\infty]\); every prime support strictly
below the closed point is retained. Nilpotents are first reduced. A
zero-dimensional local reduction has empty boundary.
For a finite reduced \(R\)-algebra \(B\), \(X_B\) is the disjoint union
of the sections at its closed centres \(\mathfrak q\) above \(\mathfrak m\),
all normalized by the same extended ideal \(\mathfrak a B_{\mathfrak q}\).
Empty centres have empty sections. Positive lower-dimensional centres
remain in this space.

Here \(C_R\) is G331's finite continuous geometric cone, and \(D_R\) its
actual decreasing completion: an element is the infimum of a nonempty
decreasing \(C_R\)-net. They are nonpositive classes, including zero;
bottom belongs to \(D_R\). G331 proves, by an independent actual-Cartier
construction, that \(C_R\) is the uniform closure of finite nonnegative
real combinations of negative unit-or-primary ideal values, and that
\(D_R\) is their stated decreasing completion. It proves full positive
descent for finite reduced injections, including all supports and
conductor-infinite fibres. Finite maxima, addition and positive scaling
preserve these classes. We use neither arbitrary supremum closure nor
flattening of iterated decreasing representations.

G345 gives the independently positive geometric mixed product
\(\mu_{R,\mathfrak a}\) of arity \(n-1\), \(n=\dim R\geq1\), on bounded
\(D_R\). It has the actual supported Cartier comparison, dense primary
ideal moments, ordinary open locality, uniform weak continuity and
bounded decreasing-net continuity with a prescribed \(D_R\) limit.
Its primary transport has exactly one measure factor. Its underlying
models are actual primary blowups with auxiliary finite normalizations,
not all proper models. Smooth specialization means precisely a smooth
affine closed germ over an algebraically closed characteristic-zero
field, plus direct DVR and empty cases. No envelope or energy theorem
is part of either input.

\begin{proposition}[Actual source model]
For any finite map \(R\to B\), and any proper primary \(I\subset R\),
the actual source model \(Y_B=\operatorname{Bl}_{IB}\operatorname{Spec}B\)
is a closed subscheme of
\(\operatorname{Spec}B\times_R\operatorname{Bl}_I\operatorname{Spec}R\).
In particular \(Y_B\to Y_R\) is finite.
Every negative actual-effective curve-nef real exceptional presentation
on \(Y_R\) pulls back to one on each nonempty source boundary component.
Its evaluation is restriction pullback on every semivaluation support.
\end{proposition}
\begin{proof}
In degree \(r\), the map \(I^r\otimes_R B\to B\) has image
\((IB)^r\). These maps give a surjection of graded algebras from the
base-changed Rees algebra to the actual source Rees algebra. Relative
Proj gives the asserted closed immersion. The base change is finite
over \(Y_R\), so the composite is finite.

The reduced source blowup has just its surviving component generic
points; on a positive-dimensional source component its generic point
does not map to \(\mathfrak m\). Otherwise that component would be finite
over the residue field, hence zero-dimensional. The blowup of the zero
ideal on a zero-dimensional component is empty. Thus a base exceptional
Cartier equation, which is a unit off the closed fibre, is nonzero at
every surviving source component generic point. On a reduced Noetherian
ring this makes it regular. Each actual effective Cartier constituent
therefore pulls back effectively; sums with nonnegative real
coefficients preserve the prescribed actual effectivity.

A proper integral source curve contracted to a closed centre maps,
by finiteness, to an integral curve in the base closed fibre.
Viewed over \(k=R/\mathfrak m\), the pulled-back bundle degree is
the old degree times the finite function-field degree. Its degree
over \(\kappa(\mathfrak q)\) divides this by the positive number
\([\kappa(\mathfrak q):k]\). Both degree conventions preserve
nonnegativity. This proves the geometric curve test.

For a source semivaluation with support \(Q\), evaluation uses its
valued field \(\operatorname{Frac}(B/Q)\), the regular generic
trivialization of the line bundle, and its centre on the primary
blowup. The base centre is obtained by composition with
\(R\to B/Q\). Pulling back the same invertible lattice gives the same
value. This proves the identity on nonminimal supports as well as on
field valuations; no map from the total fraction ring into a support
quotient was presumed.
\end{proof}

Thus core pullback is genuinely geometric before the G331 comparison.
Ideal extension, uniform limits, and actual decreasing nets then give
pullback \(C_R\to C_B\) and \(D_R\to D_B\).
If \(R\hookrightarrow B\) is injective, restriction
\(\rho:X_B\to X_R\) is a compact surjective quotient, and G331 gives
\[
 \rho^*K_R=\{U\in K_B:U\text{ is constant on every full }\rho
                       \text{ fibre}\},\qquad K=C,D.                 \tag{1}
\]
This is positive descent in the independently geometric class,
not just descent of continuous functions.

For a noninjective map put \(A=R/\ker(R\to B)\); it is reduced.
Its boundary is the closed exact image
\[
 X_A=\{v\in X_R:v(\ker(R\to B))=\infty\}.
\]
The map \(X_B\to X_A\) is surjective and satisfies (1) with \(A\)
in place of \(R\). Restriction \(K_R\to K_A\) is proved directly
by ideal restriction and the actual representing limits. We do not
assert that every \(K_A\) potential extends to \(K_R\).
Outside \(X_A\), fibres are empty. An envelope identity over the
exact image must use the cone of \(A\), not presume surjectivity
onto \(X_R\). These conventions include lower-dimensional and empty
images.

\section{Raw envelopes and the precise fibre-admissible rule}
\label{sec:envelope}
Fix \(K=C\) or \(K=D\) and a bounded finite obstacle \(f\) on \(X_R\).
Define the raw pointwise envelope
\[
 P_R^K(f)(v)=\sup\{u(v):u\in K_R,\ u\leq f\}.
                                                               \tag{2}
\]
The competitor set is nonempty: a sufficiently negative constant
\(-M=-Mv(\mathfrak a)\) belongs to \(C_R\) and lies below \(f\).
The resulting supremum is finite and nonpositive, but membership
in \(K_R\), upper semicontinuity and regularity are not asserted.
No regularization is implicit in (2).

For an injection define the fibre-admissible upstairs envelope by
the same supremum, restricting its competitors to functions in
\(K_B\) constant on every full \(\rho\)-fibre. Then
\[
 P_{B,\mathrm{fib}}^K(f\circ\rho)=\rho^*P_R^K(f),\qquad
 P_B^K(f\circ\rho)\geq\rho^*P_R^K(f).                            \tag{3}
\]
Indeed (1) is an order-preserving bijection of the two constrained
competitor sets. Evaluating that bijection at one point proves the
first formula; inclusion in the unrestricted set proves the second.
No exchange of suprema with a completion occurs.

The first envelope in (3) belongs to \(K_B\) if and only if the
downstairs envelope belongs to \(K_R\). Pullback proves one direction;
the other follows by full positive descent, because the pointwise
formula is fibre constant. Thus actual existence with positivity
transfers for this constrained operator.
If the unrestricted upstairs envelope is itself admissible, then
\[
 P_B^K(f\circ\rho)=\rho^*P_R^K(f)
 \quad\Longleftrightarrow\quad
 P_B^K(f\circ\rho)\text{ is constant on every full fibre}.        \tag{4}
\]
For sufficiency descend the actual envelope, which is then a base
competitor below \(f\); combine with (3). Necessity is immediate.
The admissibility hypothesis in (4) cannot be omitted by treating a
raw supremum as an unknown positive function.

For a noninjection the first identity in (3) holds over \(A\),
with obstacle \(f|_{X_A}\). Pullback of base \(R\) competitors
still gives \(P_B^K(f\circ\rho)\geq\rho^*P_R^K(f)\).
Equality with \(P_A^K(f|_{X_A})\) does not by itself establish
equality with the restriction of \(P_R^K(f)\).

\section{A complete continuous-obstacle normalization counterexample}
\label{sec:counter}
Let \(k=\mathbb C\) and
\[
 R=k[x,y,z]_{(x,y,z)}/(xy),\quad
 A_1=k[x,z]_{(x,z)},\quad A_2=k[y,z]_{(y,z)},\quad
 T=k[z]_{(z)}.
                                                               \tag{5}
\]
This is a two-dimensional excellent reduced local ring. Excellence
is the usual finite-type-field and localization theorem; reduction
also follows because \((xy)=(x)\cap(y)\) in the polynomial ring.
The map
\[
 R\longrightarrow S=A_1\times A_2,\qquad
 g\longmapsto(g(x,0,z),g(0,y,z))
                                                               \tag{6}
\]
identifies \(R\) with \(A_1\times_T A_2\).
To verify the assertion after localization, take a pair with equal
restriction to \(T\), subtract a lift of that common residue, and
write the remaining terms as \(xa(x,z)\) and \(yb(y,z)\).
A denominator in \(A_1\) having nonzero constant has a lift in \(R\)
with the same nonzero constant, and similarly in \(A_2\); all are
units. Thus the pair comes from \(R\). The kernel before quotient is
\((xy)\).

The algebra \(S\) is finite over \(R\). For example it is generated
as an \(R\)-module by the two idempotents \((1,0)\) and \((0,1)\):
the projections \(R\to A_i\) are surjective.
Each factor is a regular normal domain, with the fraction field of
the corresponding component of \(R\). Any element of the total
quotient ring integral over \(R\) is integral in each of these normal
factors, and hence lies in \(S\). Therefore \(S\) is the finite
normalization.

Its conductor is exactly
\[
 \mathfrak c=(R:S)=(x,y),\quad
 \mathfrak cS=xA_1\times yA_2,\quad
 R/\mathfrak c=T,\quad S/\mathfrak c=T\times T.                  \tag{7}
\]
An element of \(R\) has a common residue \(h\in T\). Multiplication
by \((1,0)\) preserves the fibre-product condition only if \(h=0\).
Conversely two zero residues remain zero after multiplication by
any pair. This proves every assertion in (7), including the exact
fibre product \(R=S\times_{S/\mathfrak c}(R/\mathfrak c)\).

Use \(\mathfrak a=\mathfrak m=(x,y,z)\).
The two clopen components of \(X_S\) are normalized by
\((x,z)\) and \((y,z)\), respectively. Let \(w_1\) kill \(x\) on
\(A_1\), let \(w_2\) kill \(y\) on \(A_2\), and on both quotients
use \(\operatorname{ord}_z\). These are centred semivaluations
with \(w_i(\mathfrak aS)=1\).
Their common restriction \(v_0\) kills \((x,y)\) and gives \(z\)
value one. Its support \((x,y)\subset R\) is nonminimal, strictly
below \(\mathfrak m\). Its conductor value is infinite.

For completeness this is exactly the whole fibre at \(v_0\).
A lift must kill the branch conductor and factor through \(T\).
A normalized centred semivaluation on the DVR \(T\) is
\(\operatorname{ord}_z\), because every nonzero element is a unit
times a power of \(z\). There is one such point on each factor.
Every other conductor-infinite base point is also \(v_0\):
it factors through this same normalized DVR.
At finite conductor value uniqueness follows directly by choosing
\(c\in\mathfrak c\) with finite value and writing
\(w(s)=v(cs)-v(c)\). Alternatively this is G331's full quotient
description. No support is removed.

On the entire \(X_R\) take
\[
 f(v)=-1-\frac1{v(y)},\qquad \frac1\infty=0.                    \tag{8}
\]
Since \(\min(v(x),v(y),v(z))=1\), \(v(y)\in[1,\infty]\).
Evaluation is continuous into the extended interval, and reciprocal
is continuous there, including infinity. Thus \(f\) is a finite
continuous obstacle with \(-2\leq f\leq-1\).
On the first branch it is identically \(-1\); on the second it is
\(-1-1/w(y)\). It is one downstairs function, agreeing at the
entire identified conductor fibre.

Let \(o_2=\operatorname{ord}_{(y,z)}\) on \(A_2\).
This is an ordinary real valuation: the associated graded ring is
\(k[y,z]\), a domain, so orders of nonzero products add.
Units have order zero; denominators in the local ring are units.
It is normalized, with both parameters of value one, and
\(f\circ\rho(o_2)=-2\).
For every \(w\in X_{A_2}\) and every proper primary ideal \(J\),
put \(d=o_2(J)\). The inclusion \(J\subset(y,z)^d\) gives
\[
 w(J)\geq d=o_2(J),                                           \tag{9}
\]
because \(w((y,z))=1\). Unit ideals satisfy (9) too.
Negation, finite nonnegative real combinations, uniform limits,
and infima of actual decreasing \(C\)-nets give
\[
 u(w)\leq u(o_2)\qquad(u\in C_{A_2}\text{ or }D_{A_2}).         \tag{10}
\]
Infinite values cause no subtraction or indeterminate operation.
If \(u\leq f\circ\rho\), (10) forces \(u\leq-2\) everywhere.
The constant \(-2\) is admissible and below the obstacle, so the
entire second-branch envelope, in either interpretation, is \(-2\).
The first-branch envelope is exactly its constant obstacle \(-1\).

These branch choices are simultaneously allowed, not merely separate
formal optima. The single semilocal boundary-power ideal
\[
 J=(x,z)\times(y,z)^2\subset S
\]
contains \((\mathfrak aS)^2\). Its negative profile is precisely
\((-1,-2)\), hence belongs to the actual ideal core and to \(C_S\).
By G331 comparison it is independently geometric. Consequently
\[
 P_S^C(f\circ\rho)=P_S^D(f\circ\rho)
       =(-1\text{ on }X_{A_1},\ -2\text{ on }X_{A_2}).          \tag{11}
\]
This entire envelope is an actual continuous admissible function.

Any downstairs competitor pulls back to an admissible second-branch
competitor. At \(v_0\), (10) therefore forces its value at most
\(-2\). The downstairs constant \(-2\) is admissible and below \(f\),
giving the exact pointwise assertion
\[
 P_R^C(f)(v_0)=P_R^D(f)(v_0)=-2.                               \tag{12}
\]
Thus, for \(K=C,D\),
\[
 P_S^K(f\circ\rho)(w_1)-P_R^K(f)(\rho w_1)=1.                 \tag{13}
\]
We do not need or claim general admissibility of \(P_R^K(f)\).
Nor do we claim it is constantly \(-2\): the primary ideal
\((x,y^2,yz,z^2)\subset R\) gives an admissible function with
value \(-1\) at the first-branch ordinary order and \(-2\) at \(v_0\).
Its first-branch values lie between \(-2\) and \(-1\), and its
second-branch value is constantly \(-2\).

The unrestricted fibre maximum of the actual function (11) fails as
well. Set \(F(v)=\max_{\rho w=v}(-1,-2)(w)\).
On second-branch monomial valuations
\[
 w_M(y)=M,\quad w_M(z)=1,\quad M\geq1,
\]
the conductor value is \(M<\infty\), the fibre is a singleton, and
\(F(\rho w_M)=-2\). These valuations converge pointwise to \(w_2\):
for a polynomial the least \(y\)-exponent dominates as \(M\) grows;
when that exponent is zero the limiting value is its \(z\)-order,
and otherwise the value tends to infinity. Local unit denominators
have value zero throughout. At the limit \(F(v_0)=-1\).
Hence \(F\) is not continuous. Its restriction to the second
branch also violates (10), since it has value \(-1\) at \(w_2\)
and \(-2\) at \(o_2\); it is not even in \(D_{A_2}\).

The algebraic reason a generic Galois argument cannot repair this
example is explicit. Normalize the target and restrict to \(A_1\):
\[
 S\otimes_R A_1=A_1\times T.
\]
The second factor maps through \(A_1/(x)\). It is a vertical
positive-dimensional component, with its own nonempty boundary,
but no generic field over \(\operatorname{Frac}(A_1)\).
Similarly the other normalized target gives \(A_2\times T\).
In full, \(S\otimes_R S=A_1\times T\times T\times A_2\),
already reduced. Conductor pullback kills each vertical quotient
factor, and its boundary point participates in the full identified
fibre. A field norm ignores that factor, while geometric descent
does not. These computations are exact fibre products of quotient
algebras, not a flat-base-change assertion.

This failure occurs in dimension two at a conductor-infinite
nonminimal support with a strictly negative continuous obstacle.
It differs from the mixed-zero literal radial obstruction on
curve discs and from an obstruction to a continuous degree trace.
It disproves unrestricted normalization envelope commutation in the
qualified geometric cone itself, while (3) remains valid.

\section{General ambient-degree finite projection for geometric products}
\label{sec:projection}
For \(n=\dim R\geq1\), define the ambient source product
\[
 \mu_B^{[R]}(\boldsymbol U)
   =\sum_{\dim B_{\mathfrak q}=n}
       f_{\mathfrak q}\,\mu_{B_{\mathfrak q},\mathfrak aB_{\mathfrak q}}
                          (\boldsymbol U|_{\mathfrak q}),\qquad
 f_{\mathfrak q}=[\kappa(\mathfrak q):\kappa(\mathfrak m)].
                                                               \tag{14}
\]
Every tuple has \(n-1\) slots. Smaller centres have zero contribution
at this ambient degree; we do not insert their different-arity
intrinsic product. Their boundary points and potentials remain in
\(X_B\). This convention defines a sum of the already geometric
G345 operators, not a new operator defined by module ranks.

For a minimal prime \(P\) of \(R\) with \(\dim R/P=n\), put
\[
 K_P=\operatorname{Frac}(R/P),\quad
 d_P(B)=\dim_{K_P}(B\otimes_R K_P),\quad
 j_P:X_{R/P}\longrightarrow X_R .
                                                               \tag{15}
\]
The map \(j_P\) evaluates through the quotient, retaining every
support containing \(P\). The degree in (15) includes inseparability.
For a noninjection it can be zero.

\begin{theorem}[Finite projection]
For every finite reduced map \(R\to B\), and every bounded tuple
\(\boldsymbol u=(u_1,\ldots,u_{n-1})\) in \(D_R\),
\[
 \rho_*\mu_B^{[R]}(\rho^*\boldsymbol u)
  =\sum_{\substack{P\in\operatorname{Min}R\\\dim R/P=n}}
       d_P(B)\,j_{P,*}\mu_{R/P,\mathfrak a(R/P)}
                                  (j_P^*\boldsymbol u).        \tag{16}
\]
The same formula with coefficient one decomposes \(\mu_R\).
If all \(d_P(B)\) equal \(d\), (16) is \(d\mu_R(\boldsymbol u)\).
\end{theorem}
\begin{proof}
For primary labels \(I_1,\ldots,I_r\), let \(H_M\) be the homogeneous
degree-\(n\) part of the eventual polynomial
\(\ell_R(M/I_1^{s_1}\cdots I_r^{s_r}M)\), for a finite \(R\)-module
\(M\). Such polynomials and their ambient convention are part of
G345's qualified algebraic kernel; their complete immediate
Swanson--Huneke 17.4.2--4 source units were freshly checked here.
We reproduce the additivity argument needed in this new application.
For an exact sequence \(0\to M_1\to M_2\to M_3\to0\),
multigraded Artin--Rees gives a fixed exponent vector \(c\) and
\[
 I^sM_1\subset I^sM_2\cap M_1\subset I^{s-c}M_1\quad(s\gg0).
 \tag{17}
\]
One proof of the needed Artin--Rees assertion is to take the graded
submodule with degree \(s\) part \(I^sM_2\cap M_1\) in the finite
multigraded Rees module of \(M_2\). Noetherian homogeneous generation
in bounded multidegrees gives the second inclusion in (17); the
first is immediate. The length of its quotient of \(M_1\) is also
the difference of the eventual quotient polynomials of \(M_2\)
and \(M_3\). The two exterior quotient polynomials in (17) have
the same top form, since a fixed exponent shift changes only
lower terms. Along every positive rational ray, dividing by the
ray parameter to power \(n\) squeezes the middle top form to the
same value. Equality on these rays is polynomial identity.
The exact quotient sequence then gives
\[
 H_{M_2}=H_{M_1}+H_{M_3}.                                    \tag{18}
\]
A module of dimension less than \(n\) contributes zero.
This argument is equality of homogeneous forms, not an invalid
coefficientwise inequality inferred from polynomial domination.

Every finite module has a finite prime filtration. Indeed an element
whose annihilator is maximal among element annihilators has prime
annihilator (apply maximality to multiplication by an element);
it gives a submodule \(R/Q\). Repeat in the quotient. Noetherianity
prevents an infinite strictly increasing sequence of the resulting
submodules, so the construction terminates. Apply (18) to a
filtration of \(B\). A factor \(R/Q\) can have dimension \(n\)
only if \(Q\) is minimal: a strict prime below it would extend an
\(n\)-chain to a longer chain in \(R\).
For top minimal \(P\), the reduced zero-dimensional local ring
\(R_P\) is the field \(K_P\). Localizing the filtration at \(P\)
leaves exactly its \(R/P\) factors, each of length one. Their
number is
\(\ell_{R_P}B_P=\dim_{K_P}(B\otimes_R K_P)=d_P(B)\).
Thus
\[
 H_B=\sum_{\dim R/P=n}d_P(B)H_{R/P}.                          \tag{19}
\]
The finite algebra in (15) is a localization of the reduced \(B\)
and is zero-dimensional over a field, hence is a product of
finite fields. Its vector-space dimension is the sum of their
full field degrees, with no separability replacement.

For any primary product \(L\), \(B/LB\) is Artinian, decomposes
over the finitely many closed centres, and has a composition
series with factors \(\kappa(\mathfrak q)\). Consequently the
exact, all-exponent identity is
\[
 \ell_R(B/LB)=\sum_{\mathfrak q}f_{\mathfrak q}
                      \ell_{B_{\mathfrak q}}(B_{\mathfrak q}/LB_{\mathfrak q}).
                                                               \tag{20}
\]
Finite integral dimension theory, or the excellent catenary
dimension formula applied componentwise, shows
\(\dim B_{\mathfrak q}\leq n\).
Comparing (19) and the degree-\(n\) part of (20) proves every
selected mixed coefficient:
\[
 \sum_{\dim B_{\mathfrak q}=n}f_{\mathfrak q}
 e_{B_{\mathfrak q}}\big((I_iB_{\mathfrak q})^{[a_i]}\big)
 =\sum_{\dim R/P=n}d_P(B)
 e_{R/P}\big((I_i(R/P))^{[a_i]}\big),\quad \sum a_i=n.          \tag{21}
\]
The factorial convention is the same on both sides. Repeated
labels, zero exponents and deletion of unused labels follow from
the whole form identity, and units have zero coefficients when
their variable occurs.

In particular take \(n-1\) separately labelled input ideals
\(I_i\), a test ideal \(J\), and divide the coefficient with all
these exponents one by the input denominators. G345's exact
geometric ideal moments turn (21) into equality of the two sides
of (16) against \(v(J)\), for rational primary profiles.
Their equality against the constant \(1=v(\mathfrak a)\) gives
equal finite mass. The rational difference lattice of these
gauges is uniformly dense in all real continuous functions on
the full \(X_R\), by G331; this includes the truncated tests
\(\min(v(g),N)=v((g)+\mathfrak a^N)\), so nonminimal supports
are tested too. Positivity and equal mass give equality of
Radon measures.

Choose uniform rational profile approximations to any \(C_R\)
tuple. Pullback and quotient restriction preserve their uniform
errors. G345's uniform weak continuity and finite mass bounds
pass (16) to that tuple. For a bounded \(D_R\) tuple, choose its
actual decreasing \(C_R\)-representations. If each limit is
at least \(-M\), all representing terms lie in \([-M,0]\).
Pullback and quotient restriction are actual decreasing
representations of known \(D\) limits. The finite product of
index sets directs the simultaneous passage. Apply G345's
prescribed-limit decreasing-net continuity to every term in
the finite sums. Continuous pushforward and these finite sums
commute with weak limits against a fixed continuous test.
This proves (16) without a measurable monotone theorem for nets
or an unproved completion membership assertion.

Apply the same argument to the module \(R\), whose top generic
lengths are one, to obtain the stated component decomposition
of \(\mu_R\). The constant-degree conclusion follows.
\end{proof}

\subsection*{The new projection is an actual Cartier projection}
We check the geometric mechanism as well as the dense moments.
Choose a common base primary model \(Y\) carrying real core
divisors \(D_1,\ldots,D_{n-1}\) and an effective exceptional
Cartier test \(E\). The actual source primary model of Section
\ref{sec:pull} maps finitely to \(Y\).
On an integral top source component, its image is an integral
top component \(Y_P\). The map is proper dominant, has equal
\(\delta\)-dimension \(n\), and has general field degree
\([L:K_P]\), which can be inseparable.
Here the base is universally catenary by excellence and
\(\delta(p)=\dim R/p\). For an immediate specialization this
drops by one: extend a saturated chain to the maximal ideal
and use catenarity. The usual finite-type extension
\(\delta(x)=\delta(p)+\operatorname{trdeg}_{\kappa(p)}\kappa(x)\)
is the one in Stacks Situation 42.7.1.

The complete immediate Stacks 42.26.3 statement and proof give
the cycle identity, before rational equivalence,
\[
 h_*\operatorname{div}(h^*s)
      =[L:K_P]\operatorname{div}(s).                          \tag{22}
\]
Its proof trivializes the line bundle locally and uses
\(\operatorname{Nm}_{L/K_P}(g)=g^{[L:K_P]}\) for a base rational
function \(g\). Its norm/divisor background is disclosed
published background; no birational degree-one premise is
substituted in (22).
The effective test equation is regular on every surviving
reduced source component. G345's actual supported Cartier
cycle decomposition retains the local thickness weights.
Only top components can supply closed-fibre primes of dimension
\(n-1\): the closed fibre of a primary blowup of a smaller
component has dimension at most its dimension minus one.
One may also see this by its effective Cartier equation and
the catenary dimension formula.

On each integral proper closed-fibre prime apply the complete
Stacks 33.45.7 numerical projection:
the intersection of \(n-1\) pulled-back line bundles is the
image intersection times its residue function-field degree,
or zero when the image dimension drops. Both schemes here
are proper over \(k=R/\mathfrak m\), not over an assumed
coefficient field inside a mixed-characteristic thickening.
G345 obtains thickening degrees by the full \(\mathfrak m\)
filtration and supported Euler polynomials; those hypotheses
are preserved here. On a source closed centre the degree over
\(k\) is \(f_{\mathfrak q}\) times the degree over
\(\kappa(\mathfrak q)\), because all Euler cohomology vector
spaces have that factor in their dimensions.
Combining (22) with numerical projection therefore gives
the general-degree supported intersection identity with
exactly the weights in (16).
Effective differences and real multilinearity handle any
Cartier test.

Thus the kernel computation identifies independently defined
positive geometric products, with actual source primary
models, regular equations and residue degrees. It does not
redefine those products by module multiplicity, or pass a
local degree through a global Chow quotient.

\subsection*{Towers, gauges, and special cases}
In a finite tower \(R\to B\to C\), localize over a top \(P\).
Write \(B\otimes_R K_P=\prod_Q L_Q\), and set
\(h_Q=\dim_{L_Q}(C\otimes_B L_Q)\). Exact vector-space
decomposition gives
\[
 d_P(C)=\sum_Q [L_Q:K_P]\,h_Q.                                \tag{23}
\]
This is the tower law for (16); source vertical components
contribute zero ambient degree. Closed residue degrees multiply
in the corresponding centre towers, so successive pushforwards
have exactly the same weights. Constant ranks \(d,e\) give \(de\).
Varying ranks retain (23), rather than a scalar averaging law.

For another base primary ideal \(\mathfrak b\), let
\(t(v)=v_{\mathfrak a}(\mathfrak b)>0\) and
\(T(v)=v/t(v)\). Each source lift has the same factor
\(t\circ\rho\). G345's transport law
\[
 \mu_{\mathfrak b}(\boldsymbol u_{\mathfrak b})
    =T_*[t\,\mu_{\mathfrak a}(\boldsymbol u_{\mathfrak a})],
 \qquad u_{\mathfrak b}(Tv)=u_{\mathfrak a}(v)/t(v),             \tag{24}
\]
commutes with every finite pushforward and component summand.
Consequently (16) is gauge covariant with one factor only;
successive factors multiply by the homogeneous gauge cocycle.

For normalization \(d_P=1\), recovering G345's full
residue-weighted conductor formula. Compatible source tuples
in \(C,D\) descend by (1) and then satisfy (16), with all
identified masses added. Lower-dimensional degree-zero
contributions do not remove their image points or assert
that conductor points have zero capacity.
For a normal domain target the degree is its generic field
degree (or the sum over dominating source fields).
For a reduced finite flat algebra of rank \(d\), every
\(B\otimes_R K_P\) has dimension \(d\), so the answer is
\(d\mu_R\). The rank of a nonreduced algebra before reduction
cannot be inserted: nilpotent generic lengths are outside
the grants.
If the exact image has dimension below \(n\), both sides of
(16) are zero ambient degree; this does not compute its
different-arity intrinsic product.
For \(n=1\), (16) has zero slots and is the actual weighted
curve reference-measure identity. For \(n=0\), every carrier
is empty and the measure is zero; no negative arity is used.

\section{Normal-domain maxima and the geometric one-slot norm trace}
\label{sec:normal}
Let \(A\) be excellent normal local domain, \(K=\operatorname{Frac}A\),
and let \(A\hookrightarrow B\) be a finite domain extension with
fraction field \(L\), \(d=[L:K]\). The source may be semilocal.
The unnormalized coefficient norm trace \(T_{B/A}\) on the full
boundary is supplied in G331, not as a product theorem:
for a boundary-power ideal \(J\subset B\), the coefficient ideal
of \(N_{L(T)/K(T)}(\sum_iT_i g_i)\) satisfies
\[
 v(\mathcal N(J))=\int_{X_B}w(J)\,d\eta_v(w),\qquad
 \eta_v(X_B)=d,\quad\operatorname{supp}\eta_v=\rho^{-1}(v).
                                                               \tag{25}
\]
Here the \(g_i\) generate \(J\) as an \(A\)-module. Its coefficients
are integral and lie in \(K\), hence lie in normal \(A\).
In a common normal field closure the norm is the product of
embedding images with generic inseparable lengths.
The polynomial Gauss minimum is multiplicative even after a
support quotient: divide by least-valued coefficients and
reduce two nonzero residue polynomials. This proves (25)
on nonminimal supports. Density of all boundary-power ideal
differences makes \(\eta_v\) independent of the chosen lift.
Every fibre point occurs under some lift to the common closure,
and then has a positive embedding weight; independence proves
full fibre support. In particular the fibres are finite.
These are the exact all-support conclusions of the qualified
grant. \(TU(v)=\int U\,d\eta_v\) preserves \(C,D\), the latter
using finite sums and actual decreasing representations.

Choose a finite normal field extension \(M/K\) containing \(L\):
take the splitting field of defining minimal polynomials,
including their inseparable roots. Its separable part is
Galois; \(G=\operatorname{Aut}_K M\) has order
\([M:K]_{\mathrm{sep}}\), and the generic inseparable length is
\([M:K]/|G|\). Each \(K\)-embedding of \(L\) into \(M\) is
the restriction of an element of \(G\), with the same number
of restrictions repeated. This follows by extending embeddings
to the algebraic closure and using that all conjugates are in
the splitting field; purely inseparable extension of an
embedding is unique.
Let \(E\) be the finite integral closure of \(A\) in \(M\).
It contains \(B\), and \(G\) preserves \(E\).
By (25) for \(E/A\), every full restriction fibre on \(X_E\)
is exactly a \(G\)-orbit, including points with nonzero support.
For \(p:X_E\to X_B\), embedding counting gives
\[
 (\rho_E)^*(T_{B/A}U)
       =\frac d{|G|}\sum_{g\in G}g^*p^*U .                   \tag{26}
\]
The formula follows first directly from the embedding sum;
it is valid for continuous and extended nonpositive functions.
No extension of an original supported semivaluation to
\(\operatorname{Frac}A\) is used.

For \(U\in C_B\) or \(D_B\), the finite maximum
\(\max_{g\in G}g^*p^*U\) belongs to the same class on \(E\),
is constant on every full fibre over \(A\), and descends by
G331 to a potential on \(A\). Its value is exactly
\[
 (\mathcal M U)(v)=\max_{\rho w=v}U(w).                        \tag{27}
\]
To see equality with the displayed maximum, extend any fibre
point of \(B\) to \(E\); every such extension is on the orbit.
Conversely every orbit point restricts to a fibre point.
Thus \(\mathcal M:K_B\to K_A\), \(K=C,D\).
For a bounded obstacle \(f\), each \(U\leq f\circ\rho\)
has \(\mathcal M U\leq f\) and \(U\leq\rho^*\mathcal M U\).
Taking competitor suprema proves the unconditional raw identity
\[
 P_B^K(f\circ\rho)=\rho^*P_A^K(f),\qquad K=C,D.                \tag{28}
\]
No actual envelope existence was required.

The same maximum conclusion holds for a finite reduced injection
over normal \(A\) if every source minimal component dominates
\(\operatorname{Spec}A\). Restrict \(U\) to each domain quotient,
apply (27), and take their finite maximum on \(A\).
Every source semivaluation kills some minimal prime, so the
union of these quotient fibres is the full fibre.
If there are vertical source components this proof does not
apply; the explicit \(A_1\times T\) example above shows why.

\begin{theorem}[One arbitrary slot over a normal domain]
Let \(n=\dim A\geq2\). For bounded \(U\in D_B\) and bounded
\(u_2,\ldots,u_{n-1}\in D_A\),
\[
 \rho_*\mu_B^{[A]}(U,\rho^*u_2,\ldots,\rho^*u_{n-1})
       =\mu_A(T_{B/A}U,u_2,\ldots,u_{n-1}).                    \tag{29}
\]
This is a new identification with the geometric G345 product.
\end{theorem}
\begin{proof}
Put \(e=[M:L]\), so \([M:K]=ed\).
The all-pulled projection (16) for \(E/B\) gives, with exact
closed residue towers,
\[
 p_*\mu_E^{[A]}(p^*U,(\rho_E)^*u_2,\ldots)
                 =e\,\mu_B^{[A]}(U,\rho^*u_2,\ldots).        \tag{30}
\]
An \(A\)-automorphism of \(E\) preserves the ambient measure
construction. On ideal tuples this follows from the
automorphism of all quotient modules and their lengths, with
the corresponding transformed ideal tests. Dense tests give
measure equality; uniform and actual bounded decreasing nets
extend it to \(D\). It also permutes closed centres preserving
their residue degrees. Since \(\rho_Eg=\rho_E\), all summands
in (26), when inserted in the selected slot and pushed to \(A\),
have the same product pushforward.
Multilinearity in that one slot therefore gives
\[
 (\rho_E)_*\mu_E^{[A]}((\rho_E)^*TU,(\rho_E)^*u_2,\ldots)
    =d\,(\rho_E)_*\mu_E^{[A]}(p^*U,(\rho_E)^*u_2,\ldots).
                                                               \tag{31}
\]
The left side is \(ed\,\mu_A(TU,u_2,\ldots)\) by (16);
the right side is \(de\,\rho_*\mu_B^{[A]}(U,\rho^*u_2,\ldots)\)
by (30). Divide by \(ed>0\).
All functions involved belong to the known bounded \(D\)
domain. This is linearity in one slot, not a product of traces.
\end{proof}

For arbitrary reduced \(B\) over normal domain \(A\), define
the dominating trace by summing the traces of \(U\) restricted
to the minimal domain components that dominate \(A\).
The component decomposition of Section \ref{sec:projection},
now applied to \(B\) and arbitrary source tuples, shows that
(29) holds with this sum trace. Other minimal components have
dimension below \(n\) and give zero ambient product.
This trace has no asserted full support on fibres containing
vertical points; those points remain relevant to (1), (3)
and (4).

\subsection*{Geometric divisor meaning and tower rules}
The trace in (29) is an independently geometric completed
potential, by G331 comparison. Its coefficient meaning is also
precise. On a normal target primary-model normalization, let
an exceptional prime \(Z\) have order \(b_Z=\operatorname{ord}_Z
(\mathfrak a)\). Its excellent DVR has finite integral closure
in \(L\), free of rank \(d\). At primes \(Z'\) above it the
uniformizer length computation gives weights
\(e_{Z'/Z}f_{Z'/Z}\), summing to \(d\), including inseparability.
For a source Cartier core \(D\) on a common source primary
refinement,
\[
 (TD)(\operatorname{ord}_Z/b_Z)
       =\frac1{b_Z}\sum_{Z'|Z}f_{Z'/Z}\operatorname{ord}_{Z'}D.
                                                               \tag{32}
\]
Here the left notation means the trace of its evaluation
function. To verify the DVR weights without a defect
assumption, the finite torsion-free integral closure over
that DVR is free; reduction modulo its uniformizer has
length equal to its rank. Artinian decomposition gives
the sum of ramification indices times residue degrees.
This agrees with (25) by field norm orders and dense ideal
tests at that base DVR.
Here the field norm order formula itself has a direct lattice
proof. For an integral nonzero \(g\), multiplication on the
finite free integral-closure module has determinant
\(N_{L/K}(g)\). Smith normal form over the DVR says its
determinant order is the length of its cokernel.
Artinian centre decomposition of that cokernel gives
\(\sum_{Z'|Z}f_{Z'/Z}\operatorname{ord}_{Z'}(g)\).
For arbitrary nonzero field \(g\), multiply by a sufficiently
large power of the base uniformizer to be integral everywhere,
then subtract the finite norm orders of that power. This
proves the formula for rational local Cartier equations too.
On a further source refinement the unique prime dominating
a given DVR has the same DVR in the same field, so its
coefficient is unchanged. Hence the right side is the
usual proper pushforward coefficient of the real divisor
on the normalized target model. It defines compatible Weil
b-divisor coefficients on the prescribed normalized primary
models. No assertion says this pushforward is Cartier on
one target model or has an actual negative-effective
Cartier presentation there. Its positive completed potential
and product identity are the assertions proved above.

In a tower of finite domain extensions over normal targets,
generic embedding sums with inseparable multiplicities give
\(T_{C/A}=T_{B/A}T_{C/B}\) whenever those traces are defined
under these hypotheses. Each embedding of the intermediate
field has the same number of extensions; their generic
Artinian lengths multiply. This proves the equality of
finite kernel sums, including all quotient supports via
(25). Alternatively it proves it on all primary gauges
and then on continuous functions by density. Finite sums
commute with actual decreasing nonpositive limits, extending
it to \(D\). The homogeneous gauge rule is
\((TU)_{\mathfrak b}(Tv)=(TU)_{\mathfrak a}(v)/t(v)\), since
every lift has the same denominator. Formula (29) then
obeys the one-factor measure transport (24).

\section{Finite flat determinant maxima and one-slot projection}
\label{sec:flat}
Assume \(R\) is excellent reduced local and \(B\) is a nonzero
reduced finite flat \(R\)-algebra of rank \(d\). It is finite free,
since \(R\) is local. This rank is the reduced rank; no
unreduced generic nilpotent multiplicity is added.
We give the extra proofs needed to identify determinant traces
and maxima with geometric operators.

\subsection*{Determinant kernels on the full boundary}
For a semilocal boundary-power ideal \(J\subset B\), choose finite
\(R\)-module generators \(g_1,\ldots,g_s\), and put
\(h=\sum T_i g_i\). Write
\[
 \det(t\,1-m_h)=t^d+\sum_{i=1}^d(-1)^i c_i(T)t^{d-i},\quad
 C_i(J)=(\text{coefficients of }c_i(T))\subset R.
                                                               \tag{33}
\]
In particular \(C_d(J)=\mathcal N_{\rm fl}(J)\) is the
determinant norm coefficient ideal.
Changing generators effects linear substitutions in \(T\);
coefficient containment in both directions proves independence.
If \(J\supset\mathfrak a^N B\), redundantly include generators
\(x\) of \(\mathfrak a^N\), as elements of \(B\).
Their pure norm coefficients are \(x^d\). Thus
\(\mathcal N_{\rm fl}(J)\) contains a power of \(\mathfrak a\)
and is unit-or-primary.

Take any \(v\in X_R\), with support \(P\), and form the valued
field \(K_v=\operatorname{Frac}(R/P)\).
The algebra \(F=B\otimes_R K_v\) has dimension \(d\), but need
not be reduced. Split its finitely many residue field
extensions in a common finite normal extension \(M/K_v\)
and extend the valuation to \(M\).
Then \(F\otimes_{K_v}M\) is a product of Artinian local
algebras with residue \(M\). On each such algebra a filtration
by powers of the maximal ideal makes multiplication by \(h\)
upper triangular, with its residue scalar on every diagonal
entry. Therefore its characteristic polynomial is
\((t-\bar h)^{l}\), where \(l\) is its length over \(M\).
The product of these expressions is (33) after extension.
This argument retains both nilpotent fibre lengths and
generic inseparability.

The residue homomorphisms \(B\to M\), followed by the extended
valuation, give centred source semivaluations \(w_j\) over \(v\).
Integrality makes them nonnegative on \(B\); the radical of
\(\mathfrak a B\) selects a closed centre and proves centring.
Polynomial Gauss valuation is multiplicative: quotient
supports first, divide each polynomial by a least-valued
coefficient, and reduce to a product of nonzero residue
polynomials with independent variables. The determinant
product now gives
\[
 v(\mathcal N_{\rm fl}(J))=\sum_j l_jw_j(J),\qquad
                      \sum_j l_j=d.                          \tag{34}
\]
All primary gauge values are finite. For a source unit-or-primary
tuple, its semilocal ideal also contains a common boundary
power, so these tests cover the dense lattice on all of \(X_B\).

Let \(\eta_v^{\rm fl}=\sum l_j\delta_{w_j}\), combining
coincident points. Its ideal moments are fixed by (34).
Its mass is \(d\). Uniform density of the G331 ideal-difference
lattice proves independence of the splitting field and
valuation lift as a finite measure.
It has full fibre support: for any \(w\) over \(v\), its
support quotient fraction field is one residue field of
\(F\). Embed it in \(M\), extend its field valuation, and
use that choice in the construction. The point \(w\) occurs
with positive length. Independence then proves that every
chosen kernel has that point in its support. Thus every
full fibre is finite, of size at most \(d\), with all
nonminimal supports retained.

Define \(\tau U(v)=\int U\,d\eta_v^{\rm fl}\).
It is positive, has continuous-function norm \(d\), and
maps continuous functions to continuous functions: (34)
proves this on dense gauges, and the fixed mass bound
passes uniform approximations.
On a rational ideal profile it is
\(-v(\mathcal N_{\rm fl}(J))/q\).
Consequently it preserves \(C_R\) by comparison and uniform
limits. Finite positive kernel sums commute with actual
decreasing nonpositive limits, including bottom. They give
actual decreasing \(C_R\)-representations, so \(\tau\)
preserves \(D\) as well.
This reconstructs the determinant valuative trace but does
not yet prove its product identification.

\subsection*{A direct flat maximum theorem}
Let \(m=\min_{\rho w=v}w(J)\), equivalently the least Gauss
value of the residue roots \(\bar h_j=\sum_i T_i\bar g_{i,j}\),
where \(\bar g_{i,j}\in M\) is the image in the \(j\)-th
residue homomorphism.
Each characteristic coefficient in (33) has Gauss value
at least \(im\). Conversely, if every coefficient had
value strictly greater than \(im\), insert a root with
value \(m\): the leading term \(t^d\) would be the unique
least-valued term of its monic equation. This is impossible.
Thus on every support
\[
 \min_{\rho w=v}w(J)=\min_{1\leq i\leq d}\frac{v(C_i(J))}{i}.
                                                               \tag{35}
\]
Zero coefficient ideals have value infinity.
The same proof handles repeated or zero roots; the last
norm coefficient is finite because \(J\) contains a boundary
power, so the minimum is finite.

The bound \(L=\sup_{X_R}v(C_d(J))/d\) is finite.
Choose an integer \(N>dL\). Replacing each \(C_i(J)\) by
\(C_i(J)+\mathfrak a^N\) leaves the minimum in (35)
unchanged, since \(N/i>L\) for every \(i\).
These replacement ideals are all allowed unit-or-primary
ideals. Hence for a source rational profile \(U=-w(J)/q\),
\[
 \mathcal M U(v):=\max_{\rho w=v}U(w)
  =\max_{1\leq i\leq d}
       \left[-\frac{v(C_i(J)+\mathfrak a^N)}{qi}\right]\in C_R.
                                                               \tag{36}
\]
The finite maximum can also be expressed as a single rational
ideal profile by taking a common denominator and an ideal sum.
Through G331 these are actual geometric potentials.

The fibre maximum is Lipschitz in uniform norm:
\(\|\mathcal M U-\mathcal M V\|\leq\|U-V\|\).
Uniform rational approximation therefore proves
\(\mathcal M(C_B)\subset C_R\).
For an actual decreasing \(C_B\)-net \(U_\alpha\downarrow U\),
each full fibre is finite. Directedness allows one index
to approximate all of its finite limiting values
simultaneously; if a value is bottom, use an arbitrary
finite lower threshold instead. This proves
\[
 \mathcal M U_\alpha\downarrow\mathcal M U .
 \]
Thus \(\mathcal M(D_B)\subset D_R\) by an actual decreasing
representation. Neither increasing closure nor regularization
was used. Applying the competitor argument of (28) proves
unrestricted raw envelope equality for every reduced finite
flat map, for \(C\) and \(D\), even when the target is
nonnormal. If one of these raw envelopes is actually
admissible, (1) and equality transfer that membership to
the other side.
This is a flat theorem with a coefficient proof, not an
extension of a Galois argument to arbitrary injections.

\begin{theorem}[Flat geometric trace slot]
For \(n=\dim R\geq2\), bounded \(U\in D_B\), and bounded
\(u_2,\ldots,u_{n-1}\in D_R\),
\[
 \rho_*\mu_B^{[R]}(U,\rho^*u_2,\ldots)
                =\mu_R(\tau U,u_2,\ldots).                   \tag{37}
\]
\end{theorem}
\begin{proof}
Let \(S\) be the finite semilocal normalization of \(R\) and
put \(B'=B\otimes_R S\). Since \(B\) is flat, \(B\to B'\)
is injective. At each normal local centre \(S_i\) of \(S\),
\(B'_i\) is finite free of rank \(d\).
Its generic algebra is
\(B\otimes_R\operatorname{Frac}(R/P_i)\), a reduced
localization of \(B\). A nilpotent in \(B'_i\) therefore
dies generically. Torsion-freeness over the domain \(S_i\)
forces it to be zero. Hence \(B'\) is already reduced,
and none of its minimal components over \(S_i\) is vertical:
localizing a finite flat algebra at a minimal prime and
using flat going-down forces its contraction to be minimal.

The determinant polynomial in (33) commutes with base change,
by its finite matrix definition. Its generator images generate
the extended ideal as an \(S_i\)-module even if tensoring the
ideal itself was not injective. Consequently on all primary
profiles
\[
 (\tau U)|_{X_{S_i}}
       =T_{B'_i/S_i}(U|_{X_{B'_i}}),                          \tag{38}
\]
where the right side is the sum of normal-domain traces
over its dominating minimal components.
Over a normal domain its generic determinant is the product
of the field norms; equality of polynomial Gauss values
gives precisely those traces, including inseparability.
Uniform limits and the finite-kernel decreasing-net passage
extend (38) to the actual bounded \(D\) tuple.

Over each \(S_i\), apply (29) and the reduced component
decomposition to get the one-slot product formula.
Push it along \(S\to R\) with the exact closed residue
weights. G345's normalization rule converts its right
side into \(\mu_R(\tau U,u_2,\ldots)\), using (38).
For its left side note that \(B\to B'\) is a finite
injection with generic degree one on every minimal
component of \(B\): at its contracted minimal \(P_i\),
normalization becomes \(\operatorname{Frac}(R/P_i)\),
and adjoining it to the source generic field changes
nothing. Formula (16) for this map, on the pulled-back
tuple \((U,\rho^*\boldsymbol u)\), converts that left
side into \(\rho_*\mu_B^{[R]}(U,\rho^*\boldsymbol u)\).
All centre residue degrees multiply in these two
pushforward routes. Lower-dimensional components of
\(R\) and \(B\) contribute zero at ambient degree \(n\)
on both routes, while their points remain present.
This proves (37).
\end{proof}

Thus a finite flat map also repairs conductor compatibility
of the trace: it is the actual downstairs completed geometric
function \(\tau U\), and its normalization restrictions are
the normal pushforward coefficients (32). This does not turn
a nonnormal Weil divisor into an actual Cartier presentation.
It is the completed geometric identification and the new
operator formula (37).

Flat towers have composing unnormalized traces. To check
the weights, specialize to the valued support field of \(R\).
The first finite free fibre algebra decomposes into
Artinian factors. A free algebra of rank \(e\) over it
has, on each residue factor after splitting, a filtration
with the same first-factor length multiplying the
second-factor residue lengths. Therefore the two-stage
kernel weights are the direct kernel weights; their
total mass is \(de\). This proves tower equality for
every source primary gauge by determinant products, and
then for every continuous test by the dense lattice.
Finite sums prove the \(D\) statement. The same calculation
is the familiar matrix identity
\(\det_R(m_g)=\operatorname{Nm}_{B/R}(\det_B(m_g))\):
it is checked after splitting each generic field factor,
and equality there is equality in reduced \(R\).
Normalized traces compose in these constant-rank towers.
Gauge rescaling gives the potential cocycle and the
one-factor measure law (24), exactly as in (29).

\subsection*{Why several traces do not give a product pushforward}
There is no nonlinear extension replacing several arbitrary
source slots by their traces. Take
\(R=k[x,y,z]_{(x,y,z)}\), \(B=R\times R\), a reduced finite
free algebra of rank two, and gauge \(\mathfrak m\).
The semilocal raw profiles \(U=(-1,0)\), \(V=(0,-1)\)
are allowed, using a primary ideal on one centre and the
unit ideal on the other. Each source-centre product
\(\mu_B(U,V)\) is zero, since one slot is zero. Its
pushforward is zero. But \(\tau U=\tau V=-1\), and
\(\mu_R(-1,-1)\) has mass \(e_R(\mathfrak m)=1\).
Thus the trace of each slot followed by multiplication
gives a nonzero measure. Dividing each trace by two
still gives positive mass \(1/4\), not zero.
This exact example preserves the validity of the
one-slot formulas and rejects the multiple-trace shortcut.

\section{Exact original scope and remaining obligations}
\label{sec:scope}
The proved package consists of actual geometric finite pullback,
full positive \(C,D\) descent for injections (using qualified
G331), the raw envelope inequality and exact fibre-admissible
identity, a complete continuous-obstacle normalization
counterexample, valid unrestricted raw identities for
normal-domain/all-dominating and reduced finite-flat maps,
the general varying-component ambient-degree product
projection, and a new one-slot geometric trace identity
at normal-domain or reduced finite-flat hypotheses.
Every new argument is given above. No ROOT348 assertion
is taken as a grant.

An unconditional demand that unrestricted envelopes commute
with normalization is false by (13). But the original
P12.6 wording does not specify that equality as its sole
target. A counterexample to that strengthened reading is
not by itself a complete negative solution of the exact
original request for geometric compatibility. We therefore
do not claim whole-original closure.

The unresolved original obligation is a general independently
geometric envelope theory with specified existence, positivity
and regularization conventions, and its full finite-map
operator calculus outside the proved raw and fibre-admissible
rules. Neither grant supplies general admissibility of
\(P_R^C(f)\) or \(P_R^D(f)\). We have not proved a geometric
pushforward for arbitrary incompatible upstairs tuples over
nonnormal nonflat targets; a generic trace cannot ignore
vertical fibres and still assert full descent. In the
noninjective case compatibility over the exact image does
not supply an extension theorem from that image to the
original ambient cone or envelope. If those are required
by a chosen complete formulation, they remain explicit
gaps, rather than being supplied by a label.

P12.3 regularity and P12.5 energy, unbounded product limits,
equation uniqueness, plurifine locality, arbitrary-model
cofinality, nilpotent generic multiplicities and literal
classical/CLD radial identifications are not imported or
proved. G345 ordinary open locality and its prescribed-limit
bounded continuity retain their exact qualifications.
All smooth statements here inherit only the exact published
BFJ scope of the two grants.
The entire original family is consequently not declared solved.

The accompanying acceptance decision and native notes set
\texttt{candidate=false}. The counterexample and scoped
positive theorems are mathematical draft claims requiring
independent review. Byte agreement, finite diagnostics and
artifact integrity establish neither mathematical correctness
nor acceptance. No compilation, external publication,
catalogue mutation, supplied program execution, subagent,
native job, private active proof or hot feedback was used.
